\documentclass[a4paper,10pt]{article}
\usepackage[francais]{babel}
\usepackage{url}

\selectlanguage{francais}
\input{letter-ensiie.sty}

\selectlanguage{francais}
\enCouleur

\begin{document}


\name{St\'ephane \textsc{Genaud}}
\fonction{Professor at ENSIIE}
\telephoneillkirch{+33 (0)3688 54542}
\email{genaud@ensiie.fr}

\begin{letter}[Recommendation de Nestor \textsc{Demeure}]{à}%
{Ecole doctorale de mathématiques Hadamard}


Chers Collègues,

Je  connais Nestor Demeure  depuis sa  première année  de  scolarité dans  l'école
d'ingénieur  en  informatique   ENSIIE,  dont  je  dirige   l'antenne  locale  à
Strasbourg, et dans laquelle j'interviens également comme enseignant et y
supervise les stages.
Je donne des éléments d'appréciation de la candidature de Nestor à ces deux titres.

\textbf{Résultats Scolaires}

La formation d'ingénieur de l'ENSIIE est dispensée à la fois sur le site
historique principal, à Evry, et à son antenne de Strasbourg selon les modalités
suivantes: un tronc commun se déroule sur deux semestres, avec des contenus 
pédagogiques et des examens communs aux deux sites, puis les semestres 3, 4 et 5
sont constitués à 80\% d'options axées sur les spécialités recherche des sites.
Le dernier semestre est un stage long de type professionnel ou recherche. 

Je n'ai pu obtenir dans le temps imparti les résultats de Nestor que par rapport
à la promotion Strasbourgeoise, soit environ  50 élèves. Le tableau suivant fait
apparaître  les résultats généraux de Nestor. Je n'ai aucun doute que ramené à
l'ensemble de la promotion de l'école, Nestor s'est également classé dans les
tous premiers de sa promotion.

\begin{center}
\begin{tabular}{ccc}
semestre & rang/effectif & moyenne générale\\
\hline \\
S1 & 2/51 & 15,73/20\\
S2 & 2/51 & 16,07/20\\
S3 & 3/38 & 15,26/20\\
S4 & 2/43 & 16,82/20\\
\hline \\
\end{tabular}
\end{center}

Ses résultats scolaires sont  tout à fait remarquables, et sont  le reflet de la
personnalité de  Nestor. Attentif, discret, travailleur,  extrêmement rigoureux,
Nestor  excelle  dans  tous  les   domaines,  en  particulier  mathématiques  et
informatique.  D'un point  de vue  scolaire, je  placerais Nestor  dans les  10\%
meilleurs étudiants que j'ai cottoyé dans ma carrière.
  

\textbf{Motivation pour la Recherche}

Nestor  a choisi  comme stage  de deuxième  année, un  stage orienté  recherche,
intitulé \emph{Biostatistical  analysis of  large-scale} dans le  laboratoire de
biologie moléculaire IGBMC  (UMR Inserm, CNRS et Université  de Strasbourg).  Le
chercheur encadrant le  stage lui a demandé de revoir  le pipeline de traitement
des    données    expérimentales    pour   en    améliorer    l'efficacité    et
l'exploitation.  Lors de  sa  soutenance  de stage,  l'encadrant  a souligné  la
qualité exceptionnelle du travail, apportant de nombreuses idées aux pratiques
actuelles d'analyse des données de l'équipe.

Personnellement, je  suis entièrement  convaincu que  Nestor possède  toutes les
qualités pour mener une carrière prometteuse  dans le monde de la recherche.  Sa
curiosité, sa  large culture scientifique  et son  travail posé sont  des atouts
incontestables  pour un  doctorat.  La  thèse envisagée  sur la  question de  la
\textit{Gestion du  compromis entre la  performance et  la précision de  code de
  calcul} au sein du laboratoire LRC MESO  me semble tout à fait dans le domaine
de compétences de Nestor.

Pour ces  raisons, je  recommande au  plus haut point  la candidature  de Nestor
Demeure pour mener ce travail de recherche.


Je vous prie, Madame, Monsieur, d'agréer mes sentiments respectueux.

\end{letter}
\begin{flushright}
\includegraphics[width=.2\textwidth]{signgenaud.jpg}
\end{flushright}
\end{document}
